% Vietnamese translation for OI Public Library
% Source-PDF: IOI中国国家候选队论文/2000/骆骥论文.pdf
\documentclass[11pt]{article}
\usepackage{oipl-vn}

\newcommand{\Oh}{\mathcal{O}}

\title{Xây Dựng Và Lựa Chọn Mô Hình Toán Học}
\author{Luo Ji\\Trường Trung học số 1 Vu Hồ, An Huy}
\date{2000}

\begin{document}
\maketitle

\origtitle{Establishment and selection of mathematical models}
\sourcepath{IOI Candidate Team Papers/2000/Luo Ji.pdf}

\begin{abstract}
Bài viết bàn về khâu then chốt trong giải bài thi tin học: xây dựng và lựa
chọn mô hình toán học. Tác giả phân tích tầm quan trọng của việc đi từ nguyên
mẫu thông tin tới mô hình toán học, nêu quá trình ``thực tế -- lý thuyết --
thực tế'', rồi chia phương pháp xây dựng mô hình thành phân tích cơ chế và
phân tích thống kê. Sau đó bài viết thảo luận các tiêu chuẩn lựa chọn mô hình
và tổng kết một quy trình chung.
\end{abstract}

\noindent\textbf{Từ khóa:} nguyên mẫu thông tin; mô hình toán học; xây dựng mô
hình.

\section{Từ nguyên mẫu thông tin tới mô hình toán học}

Trong thế giới rộng lớn, các sự vật ta đối diện rất đa dạng và phức tạp. Chúng
đều được tạo bởi nhiều loại thông tin. Những mô tả bằng ngôn ngữ tự nhiên về
bề mặt của một vấn đề khách quan trong thế giới thực được gọi là
\term{nguyên mẫu thông tin}. Trong thi tin học, bài toán ta gặp cũng là các
nguyên mẫu thông tin.

Nguyên mẫu thông tin thường bị bao phủ bởi nhiều chi tiết rối rắm, che lấp
những thuộc tính quan trọng. Vì vậy ta khó trực tiếp tìm đáp án từ nguyên mẫu.
Cần có một cách ``cải tạo'' nguyên mẫu, giữ lại thuộc tính quan trọng nhưng
làm cho nó có thể nghiên cứu được.

Ta chuyển các thuộc tính của nguyên mẫu thông tin vào một mô hình. Mô hình là
sự mô phỏng các thuộc tính của vấn đề khách quan; nó đại diện cho nguyên mẫu
và là đối tượng ta có thể nghiên cứu. Nếu dùng ngôn ngữ toán học để trừu tượng
hóa, phiên dịch và quy nạp nguyên mẫu thông tin, ta thu được một
\term{mô hình toán học}. Mô hình ấy vừa giữ thuộc tính của bài toán gốc, vừa
có khả năng được nghiên cứu bằng toán.

Quá trình giải bài trong thi tin học có thể xem là:
\[
  \text{thực tế}\longrightarrow \text{lý thuyết}\longrightarrow \text{thực tế}.
\]
Nguyên mẫu thông tin là vấn đề thực tế; mô hình toán học là mô hình lý thuyết;
nghiên cứu mô hình đem lại lời giải cho vấn đề thực tế. Bản gốc minh họa điều
này bằng sơ đồ: nguyên mẫu thông tin được xây dựng mô hình thành mô hình toán
học, mô hình dẫn tới quyết sách thuật toán, và thuật toán cho lời giải của
nguyên mẫu.

Do đó, trong toàn bộ quá trình phân tích và giải bài, bước chuyển từ nguyên
mẫu thông tin sang mô hình toán học là cực kỳ quan trọng. Bước ấy được gọi là
\term{xây dựng mô hình}.

\section{Xây dựng mô hình toán học}

Từ kinh nghiệm thực tế, xây dựng mô hình không có khuôn mẫu cố định và phương
pháp rất đa dạng. Nhưng nhìn chung có thể chia thành hai nhóm lớn:
\term{phân tích cơ chế} và \term{phân tích thống kê}.

\subsection{Phân tích cơ chế}

Phân tích cơ chế là phương pháp dựa trên đặc tính của sự vật khách quan, phân
tích cơ chế bên trong và làm rõ các quan hệ, rồi trong điều kiện trừu tượng
thích hợp dùng công cụ toán học để thu được mô hình mô tả thuộc tính của sự
vật.

Trong thi tin học, ta thường dùng phương pháp này để xây dựng mô hình, sau đó
chuyển mô hình thành thuật toán. Sơ đồ của bản gốc có thể diễn giải như sau:
từ nguyên mẫu thông tin, ta dùng phân tích cơ chế để trừu tượng hóa thành mô
hình toán học; từ mô hình, ta tìm tương ứng lý thuyết tới thuật toán; thực thi
thuật toán để nhận được lời giải.

Các phương pháp phân tích cơ chế có nhiều dạng. Trong thực tế thi đấu, ta
thường dùng pha trộn, nên việc phân loại không tuyệt đối. Theo các tầng mức
khác nhau của xây dựng mô hình, bài viết chia thành bốn loại.

\subsubsection{Xây dựng mô hình trực tiếp}

Khi nguyên mẫu thông tin tương đối đơn giản và thuộc tính của nó rõ ràng, ta
thường dùng phương pháp xây dựng mô hình trực tiếp. Có thể hiểu đây là việc
tổng hợp các thuộc tính hiển nhiên của nguyên mẫu theo phương pháp toán học để
thu được mô hình. Mô hình này có tính nhắm đích mạnh, nhưng thường không có ý
nghĩa phổ quát.

Thực chất, xây dựng mô hình trực tiếp cũng là một dạng sáng tạo, nhưng tư duy
tương đối đơn giản và không có phát hiện lý thuyết mới; tác giả gọi đó là
\term{sáng tạo cấp một}. Ví dụ, bài xếp hạng trong NOI 1997 có các thuộc tính
rất rõ ràng, bản thân đề bài đã có mức trừu tượng nhất định và nhiều tính chất
được mô tả gần bằng ngôn ngữ toán học, nên có thể trực tiếp xây dựng mô hình
đơn giản để giải.

\subsubsection{Áp dụng mô hình quen thuộc}

Khi xây dựng mô hình, ta thường hướng việc trừu tượng hóa về một mô hình quen
thuộc đã biết. Nếu khớp, ta lập tức dùng mô hình ấy và áp dụng nguyên thuật
toán tương ứng.

\paragraph{Ví dụ: xâu 01, NOI 1999.}
Đề bài yêu cầu thiết kế một xâu nhị phân
\[
  S=s_1s_2\cdots s_N
\]
thỏa:
\begin{enumerate}
  \item \(s_i=0\) hoặc \(s_i=1\), với \(1\le i\le N\);
  \item với mọi đoạn liên tiếp độ dài \(L_0\), số ký tự \(0\) nằm giữa
  \(A_0\) và \(B_0\);
  \item với mọi đoạn liên tiếp độ dài \(L_1\), số ký tự \(1\) nằm giữa
  \(A_1\) và \(B_1\).
\end{enumerate}

Mô hình toán học bị che khá kỹ trong nguyên mẫu thông tin. Trước hết dùng kỹ
thuật tổng tiền tố. Từ xâu \(S\), xây dựng dãy \(T\):
\[
T_i=
\begin{cases}
0, & i=0,\\
\displaystyle\sum_{j=1}^{i}s_j, & i>0.
\end{cases}
\]
Hai dãy \(S\) và \(T\) tương ứng một-một. Các điều kiện trên chuyển thành:
\[
  0\le T_i-T_{i-1}\le 1 \qquad (1\le i\le N),
\]
\[
  A_0\le L_0-(T_{j+L_0}-T_j)\le B_0 \qquad (0\le j\le N-L_0),
\]
\[
  A_1\le T_{j+L_1}-T_j\le B_1 \qquad (0\le j\le N-L_1).
\]
Cách mô tả bằng toán đã trừu tượng hơn nhiều so với ngôn ngữ tự nhiên, nhưng
ta vẫn mới hoàn thành bước đầu. Bài toán trở thành giải một hệ bất đẳng thức.

Điều này gợi tới bài ``quy hoạch công trình'' trong vòng chọn đội Trung Quốc
cho IOI 1997, tức mô hình đường đi ngắn nhất với ràng buộc sai phân. Mỗi bất
đẳng thức chỉ liên quan hai biến, nên có thể xem hai biến là hai đỉnh và bất
đẳng thức là cạnh biểu diễn quan hệ giữa chúng. Từ đó dựng đồ thị; đặt
\(T_0=0\), giá trị \(T_i\) là độ dài đường đi ngắn nhất từ \(V_0\) tới \(V_i\).
Vì có cạnh âm, dùng phép lặp thư giãn. Như vậy mô hình đồ thị được thiết lập
và thuật toán hiệu quả xuất hiện.

Tác giả gọi cách này là \term{bắt chước cấp một}: liên tưởng tới mô hình quen
thuộc và áp dụng nó. Cách làm rất thường dùng, nhưng có thể giới hạn khả năng
sáng tạo và đôi khi hơi máy móc.

\subsubsection{Sửa đổi mô hình quen thuộc theo đặc thù bài toán}

Một nguyên mẫu thông tin luôn có thuộc tính riêng. Vì vậy ta nên sửa đổi mô
hình quen thuộc sao cho đưa được thuộc tính riêng ấy vào, nhằm tối ưu mô hình
và thuật toán.

\paragraph{Ví dụ: tìm số lớn thứ \(k\).}
Cho \(n\) số phân biệt, hãy tìm số lớn thứ \(k\), với
\[
  1\le k\le n\le 10000.
\]
Cách trực tiếp là áp dụng mô hình sắp xếp nhanh: sắp xếp toàn bộ dãy rồi lấy
số lớn thứ \(k\), độ phức tạp \(\Oh(n\log n)\).

Nhưng thuộc tính đặc thù của nguyên mẫu là ta chỉ cần số lớn thứ \(k\), không
cần toàn bộ dãy có thứ tự. Trong quicksort, khi chọn được một điểm chia \(x\),
dãy tách thành hai phần. Nếu vị trí của \(x\) lớn hơn \(k\), đáp án nằm ở
phần trái và không cần xử lý phần phải; nếu nhỏ hơn \(k\), đáp án nằm ở phần
phải; nếu bằng \(k\), ta đã tìm được đáp án. Vì vậy mỗi lần chỉ cần đệ quy
vào nhiều nhất một phía.

Nếu điểm chia thường gần trung vị, có thể dùng ngẫu nhiên hóa để bảo đảm xấp
xỉ, chi phí trung bình là
\[
  n+\frac n2+\frac n4+\frac n8+\cdots = 2n,
\]
tức \(\Oh(n)\), thấp hơn \(\Oh(n\log n)\).

Đây là \term{bắt chước cấp hai}: vẫn bắt nguồn từ mô hình quen thuộc, nhưng
có sửa đổi cục bộ theo thuộc tính riêng của bài toán.

\subsubsection{Sáng tạo tổng hợp}

Nhiều bài toán khó giải bằng bắt chước, trong khi thuộc tính nguyên mẫu cũng
không rõ để xây dựng trực tiếp. Khi đó cần sáng tạo tổng hợp. Phương pháp này
dựa vào tri thức toán học, dùng mô hình hoặc phương pháp đã biết để phân tích
thuộc tính của nguyên mẫu, rồi sáng tạo một mô hình hoặc phương pháp mới có
tính ứng dụng rộng hơn.

\paragraph{Ví dụ: vệ tinh phủ, NOI 1997.}
Bài \textit{Cover} là một bài tổng hợp tốt, kiểm tra năng lực tưởng tượng
không gian, khái niệm toán học và kỹ năng lập trình. Tác giả tập trung vào
đoạn suy nghĩ dẫn tới mô hình rời rạc hóa.

\textit{Cover} là một bài thống kê dữ liệu trong hệ tọa độ tĩnh xuất hiện khá
sớm. Những bài cùng loại về sau thường có dữ liệu lớn và số chiều cao. Khi
phân tích, ta mong muốn chia dữ liệu quy mô lớn theo một nguyên tắc nào đó để
trị từng phần. Đây là liên tưởng từ tư tưởng chia để trị, với mục tiêu hạ số
chiều.

Khó khăn là chia như thế nào. Nếu chia đều các phần tử trong hệ tọa độ tĩnh,
ta nhận được một lưới đều không phụ thuộc dữ liệu bài toán. Cách chia cố định
như vậy mù quáng. Với bài \textit{Cover}, nếu hai hình chữ nhật phủ nhau mà ta
vẫn chia đều theo các đường không liên quan, nhiều phần chia không hề làm giảm
độ phức tạp tư duy, thậm chí còn tăng chi phí.

Vì vậy ta kỳ vọng một cách chia phụ thuộc vào chính dữ liệu. Với hai hình chữ
nhật, quan hệ theo hoành độ thể hiện ở những đường qua cạnh của chúng; chỉ cần
chia theo các đường đó. Các tọa độ dùng để chia được gọi là \term{điểm rời
rạc}. Dựa vào điểm rời rạc của các hình, ta xây dựng mô hình rời rạc hóa, sao
cho các bài toán con không còn quan hệ tô-pô phức tạp. Bản chất của mô hình
này là chia bài toán thành các phần ``càng lớn càng tốt'', số lượng ít nhưng
quan hệ đơn giản và hiệu quả cao.

Mô hình rời rạc hóa của \textit{Cover} là kết quả của sáng tạo tổng hợp. Nó có
khái niệm mới, tư tưởng mới và phương pháp mới; tác giả gọi đây là
\term{sáng tạo cấp hai}.

\subsection{Phân tích thống kê}

Phân tích thống kê được dùng khi ta tạm thời không tìm được cơ chế đặc trưng
của sự vật. Khi đó, thông qua thử nghiệm hoặc một chương trình đơn giản, ta
thu được một số dữ liệu, tức một phần lời giải của bài toán; rồi dùng tri thức
thống kê toán học để xử lý dữ liệu và rút ra mô hình toán học.

So với phân tích cơ chế, đây là tư duy ngược. Trong thi tin học, ta thường
dùng tay hoặc chương trình đơn giản để tìm một số lời giải nhỏ, sau đó phân
tích quy luật trong đó.

\paragraph{Ví dụ: bài cực trị, NOI 1995.}
Cho số nguyên \(m,n\) thỏa:
\[
  m,n\in\{1,2,\ldots,k\},\qquad 1\le k\le 10^9,
\]
\[
  (n^2-mn-m^2)^2=1.
\]
Nhập \(k\), tìm một cặp \(m,n\) thỏa điều kiện và làm \(m^2+n^2\) lớn nhất.

Đây là bài có tính toán học mạnh. Trước biểu thức quan hệ giữa \(n\) và \(m\),
ta khó tiếp tục trừu tượng hóa bằng phân tích cơ chế. Lúc này phân tích thống
kê phát huy tác dụng. Với vài giá trị nhỏ:
\[
\begin{array}{c|rrrrrrr}
k & 2&3&4&8&16&32&64\\ \hline
N & 2&3&3&8&13&21&55\\
M & 1&2&2&5&8&13&34
\end{array}
\]
Ta thấy \(N,M\) tăng theo dạng dãy Fibonacci. Vì vậy phỏng đoán: lời giải là
hai số Fibonacci kề nhau lớn nhất không vượt quá \(k\).

Nếu \((n,m)\) là một nghiệm của
\[
  (n^2-mn-m^2)^2=1,
\]
thì theo quan hệ Fibonacci, \((n+m,n)\) cũng phải là nghiệm. Kiểm tra:
\[
\begin{aligned}
&\bigl((n+m)^2-(m+n)n-n^2\bigr)^2\\
&\quad=(n^2+mn+n^2-n^2-2mn-m^2)^2\\
&\quad=(n^2-mn-m^2)^2=1.
\end{aligned}
\]
Các bước có thể đảo ngược, nên phỏng đoán được củng cố. Bài toán ban đầu được
giải bằng cách sinh Fibonacci.

Phân tích thống kê rất hữu dụng trong ví dụ này. Nhưng không phải nguyên mẫu
nào cũng dùng được phương pháp này: đôi khi ta không thể tìm được lời giải bộ
phận, hoặc không thể từ lời giải nhỏ rút ra quy luật.

\section{Lựa chọn mô hình toán học}

Một nguyên mẫu thông tin có thể tương ứng với nhiều mô hình toán học. Khi đó
ta phải lựa chọn. Vì trong thi đấu mô hình cuối cùng phải chuyển thành thuật
toán, đánh giá mô hình chủ yếu là đánh giá thuật toán tương ứng.

Các nguyên tắc đánh giá thường là:
\begin{enumerate}
  \item \term{Độ phức tạp thời gian.} Một thuật toán tốt thường có hiệu suất
  cao. Đây là tiêu chuẩn quan trọng nhất vì bài thi có giới hạn thời gian.
  \item \term{Độ phức tạp không gian.} Do giới hạn bộ nhớ, mô hình phải dẫn
  tới thuật toán không vượt quá bộ nhớ. Không nhất thiết càng ít bộ nhớ càng
  tốt; chỉ cần nằm trong giới hạn.
  \item \term{Độ phức tạp lập trình.} Nếu thuật toán quá rắc rối để hiện thực
  trong cuộc thi, nó cũng không thuận lợi.
\end{enumerate}

Khi có nhiều mô hình, thường có hai tình huống. Thứ nhất là cùng một tư tưởng
nhưng nhiều cách hiện thực khác nhau; ví dụ cùng là quy hoạch động nhưng có
hai cách hiện thực khác nhau. Thứ hai là các tư tưởng hoàn toàn khác nhau; ví
dụ một bài có thể dùng quy hoạch động hoặc luồng cực đại. Trong nhiều trường
hợp, nếu cả hai đều có thể dùng, ta thường chọn quy hoạch động vì mô hình hóa
luồng cực đại có thể khó hơn: phải xác định ý nghĩa của đỉnh, cạnh, lưu lượng,
chi phí, v.v.

Ngoài ba tiêu chuẩn trên, trong quá trình phân tích còn có \term{độ phức tạp
tư duy}: thời gian và sức lực cần để đi từ nguyên mẫu tới mô hình. Trong thi
đấu, ta hiếm khi xây nhiều mô hình hoàn chỉnh rồi mới chọn, mà thường vừa
phân tích vừa lựa chọn hướng tốt hơn.

\section{Tổng kết}

Từ các phân tích trên, có thể xem việc giải một bài toán là tìm con đường tốt
nhất từ ``nguyên mẫu thông tin'' tới ``toàn bộ lời giải''. Xây dựng mô hình
bằng phân tích cơ chế thường gồm các bước:
\begin{enumerate}
  \item \term{Diễn đạt lại bài toán bằng lời của mình.} Đây là chuyển hóa
  nguyên mẫu ở tầng ngôn ngữ tự nhiên, giúp nắm thuộc tính chính.
  \item \term{Diễn đạt lại bằng ngôn ngữ toán học.} Nắm thuộc tính chính rồi
  dùng ngôn ngữ toán học mô tả, tức trừu tượng hóa nguyên mẫu.
  \item \term{Liên tưởng tri thức liên quan.} Nếu mô tả toán học đã đơn giản
  rõ ràng, có thể trực tiếp lập thuật toán hoặc mô hình đơn giản. Nếu không,
  liên tưởng tới nguyên mẫu tương tự hoặc mô hình đã biết.
  \item \term{Rút ra sự thật hữu ích.} Nếu mô hình liên tưởng được có ích, có
  thể áp dụng; xa hơn, dựa trên thuộc tính riêng của nguyên mẫu để cải tiến.
  \item \term{Sáng tạo mạnh dạn.} Sau các bước trên, trong đầu có thể xuất
  hiện một mô hình mới còn mơ hồ. Đừng bỏ qua nó; sáng tạo có thể dẫn tới một
  mô hình hoặc hệ phương pháp có phạm vi áp dụng rộng.
\end{enumerate}

Phân tích thống kê cũng quan trọng. Thông thường, ta trước hết dùng phân tích
cơ chế; nếu không tiến triển, hoặc bài toán có tính toán học mạnh và dễ tìm
lời giải nhỏ, ta cân nhắc phân tích thống kê. Hai phương pháp không tách rời:
kết luận từ phân tích cơ chế có thể giúp thống kê, và quy luật thống kê cũng
có thể quay lại phục vụ phân tích cơ chế.

Sau khi lập mô hình, công việc mới hoàn thành một nửa. Nếu có nhiều mô hình,
cần xét thời gian, không gian và độ phức tạp lập trình để lựa chọn. Cuối cùng,
ta còn phải liên tục hoàn thiện mô hình để thu được một mô hình chính xác và
hiệu quả.

\section*{Phụ lục}

\begin{itemize}
  \item Bài gốc kèm phát biểu bài \textit{Xâu 01}: cho các số
  \(N,A_0,B_0,L_0,A_1,B_1,L_1\), yêu cầu dựng xâu nhị phân thỏa các ràng buộc
  số lượng \(0\) và \(1\) trong mọi đoạn liên tiếp.
  \item Phụ lục chương trình gồm ba ví dụ: \textit{Xâu 01} của NOI 1999, tìm
  số lớn thứ \(k\), và bài cực trị của NOI 1995. Trong bản dịch này, nội dung
  thuật toán của ba chương trình đã được diễn giải trong các mục tương ứng
  thay vì chép lại mã Pascal dài có chú thích tiếng Trung.
\end{itemize}

\section*{Tài liệu tham khảo}

Bản PDF chỉ trích dẫn tài liệu qua các chú thích trong bài; tiêu điểm của bài
là phương pháp xây dựng và lựa chọn mô hình trong bối cảnh thi tin học.

\end{document}
